% ============================================================================
%  chapter4.tex
%  Data Structures and Algorithms in C (The Mega Chapter)
% ============================================================================

\chapter{Data Structures and Algorithms in C}

\section{Welcome to the Menagerie}

Everything in the previous chapters was preparation for a single
realization: \emph{a data structure is a disciplined web of pointers}. A
linked list is a struct pointing at a struct. A tree is a struct with
two or more pointers to structs. A graph is an array of lists of
structs. A hash table is an array of pointers to structs, chosen by
hashing. There is nothing else under the hood---only structs, pointers,
and the ownership rules we insisted upon.

This chapter is the biggest in the book, and deliberately so. We
implement, draw, and dissect the entire family of structures a systems
programmer reaches for daily. For each one we answer the same four
questions, the ones that settle every design:

\begin{verbatim}
   1. SHAPE  : what does the struct look like?       (draw it)
   2. COST   : what do search/insert/delete cost?    (O(...))
   3. OWNER  : who allocates, who frees?             (write it down)
   4. CACHE  : how do the bytes sit in memory?       (think about it)
\end{verbatim}

The code follows the house style without exception: pointers move on
their own lines, \code{size\_t} is reserved for loop counters, every
\code{malloc} is checked, and every diagram is drawn before any code is
trusted.

\section{Linked Lists}

\subsection{The Singly Linked List}

The singly linked list is the first structure, the seed from which the
others grow. A node holds a value and one pointer: the address of the
next node.

\begin{lstlisting}[caption={The singly linked node.}]
typedef struct node {
    int value;
    struct node *next;
} node;
\end{lstlisting}

\begin{verbatim}
   head
    |
    v
   +-------+      +-------+      +-------+
   | val=3 |      | val=7 |      | val=1 |
   | next -----> | next -----> | next -----> NULL
   +-------+      +-------+      +-------+

   - The list is named by the single pointer "head".
   - Every "next" is an arrow to the following node.
   - The final "next" is NULL: the arrow that points at nothing.
   - You may only walk FORWARD; to reach node k you must pass through
     nodes 0..k-1, one dereference at a time.
\end{verbatim}

Push at the front is the operation we already know from Chapter~3: it
rewrites the caller's \code{head} through a pointer-to-pointer.

\begin{lstlisting}[caption={Push at the front. Returns 1 on success.}]
int list_push_front(node **head, int value) {
    node *new_node;

    new_node = malloc(sizeof(node));
    if (new_node == NULL) {
        return 0;
    }
    new_node->value = value;
    new_node->next = *head;
    *head = new_node;
    return 1;
}
\end{lstlisting}

Singly lists reward one trick above all others: keeping the address of
the \emph{link} you are interested in, rather than the node. Inserting
in sorted order and deleting both collapse to the same pattern---walk a
\code{node **} and splice.

\begin{lstlisting}[caption={Insert keeping the list sorted ascending.}]
int list_insert_sorted(node **head, int value) {
    node *new_node;
    node **cur;

    new_node = malloc(sizeof(node));
    if (new_node == NULL) {
        return 0;
    }
    new_node->value = value;
    new_node->next = NULL;

    /* cur is the ADDRESS of a "next" link, not the address of a node.
       It walks until the first node whose value is >= value.        */
    cur = head;
    while (*cur != NULL && (*cur)->value < value) {
        cur = &(*cur)->next;
    }

    /* Splice: the new node points at the node cur names, and the
       link cur names now points at the new node.                  */
    new_node->next = *cur;
    *cur = new_node;

    return 1;
}
\end{lstlisting}

\begin{verbatim}
   Insert 5 into  3 -> 7 -> NULL

   head
    |
    v
   +-------+      +-------+
   | val=3 |      | val=7 |
   | next -----> | next -----> NULL
   +-------+      +-------+

   cur begins as &head (an address OF A POINTER).
     step 1: *cur = 3-node, 3 < 5, so cur = &(3-node.next)
     step 2: *cur = 7-node, 7 not < 5, stop

   Now:
     new_node->next = *cur;   (new node -> 7-node)
     *cur = new_node;         (3-node.next -> new node)

   head
    |
    v
   +-------+      +-------+      +-------+
   | val=3 |      | val=5 |      | val=7 |
   | next -----> | next -----> | next -----> NULL
   +-------+      +-------+      +-------+

   Removing a node uses the same trick, and never needs a special case
   for "removing the head".
\end{verbatim}

\begin{lstlisting}[caption={Remove the first node holding value.}]
int list_remove(node **head, int value) {
    node **cur;
    node *victim;

    cur = head;
    while (*cur != NULL && (*cur)->value != value) {
        cur = &(*cur)->next;
    }

    if (*cur == NULL) {
        return 0;              /* not present */
    }

    victim = *cur;
    *cur = victim->next;       /* the predecessor now skips victim */
    free(victim);
    return 1;
}
\end{lstlisting}

Reading \code{list\_remove}: because \code{cur} is an address of a
link, unlinking is one assignment. If the victim was the head,
\code{cur} is \code{\&head} and \code{*cur} updates the caller's head
directly. No branch, no special case, no bug.

\begin{verbatim}
   Removing the 7-node from  3 -> 7 -> 9 -> NULL:

   BEFORE                                  AFTER
   +-------+   +-------+   +-------+       +-------+   +-------+
   | val=3 |   | val=7 |   | val=9 |       | val=3 |   | val=9 |
   | next ----->| next ----->| next --->   | next ----->| next --->
   +-------+   +-------+   +-------+       +-------+   +-------+
                 ^ victim

   *cur (which is 3-node.next) is set to victim->next (the 9-node).
   Then free(victim). The 7-node is gone; the chain is unbroken.
\end{verbatim}

Traversal and destruction are the loops we saw in Chapter~3: walk with
\code{cur = cur->next}, and in \code{free\_all} save the successor
before freeing. Search is the same walk returning the matching node, or
\code{NULL}.

\begin{insight}[Why we keep the link address, not the node]
The singly linked list has only forward arrows. To delete a node you
must be able to reach the arrow that points at it---the \emph{previous}
node's \code{next}. Rather than tracking the previous node in a second
variable, we walk a pointer \emph{to} the link itself. \code{node **}
is not advanced; it is exactly the right precision: it is the address
of the arrow.
\end{insight}

\subsection{The Doubly Linked List}

Add a second pointer to each node and you can walk both ways, delete a
node in \code{O(1)} given its address, and insert before a node without
searching. The cost is one extra pointer per node.

\begin{lstlisting}[caption={The doubly linked node.}]
typedef struct dnode {
    int value;
    struct dnode *prev;
    struct dnode *next;
} dnode;
\end{lstlisting}

\begin{verbatim}
   head
    |
    v
   +--------+      +--------+      +--------+
   | prev=NULL |   | prev=3 |       | prev=7 |
   |  value=3 |   | value=7 |       | value=9 |
   | next=7 ----+  | next=9 ----+   | next=NULL |
   +--------+  |  +--------+  |    +--------+
        ^       +----->|       +----->|
        |               |               |
        +----(the prev arrows)----------+

   Two chains of arrows: "next" runs forward, "prev" runs backward.
   head->next is the first node; head->prev is NULL (or the tail, in
   the circular variant below).
\end{verbatim}

Inserting after a given node touches four arrows. Do them in an order
that never reads a node you are about to unlink.

\begin{lstlisting}[caption={Insert a new node right after node at.}]
int dlist_insert_after(dnode *at, int value) {
    dnode *new_node;

    new_node = malloc(sizeof(dnode));
    if (new_node == NULL) {
        return 0;
    }
    new_node->value = value;

    /* Wire the new node into the gap. */
    new_node->prev = at;
    new_node->next = at->next;

    /* Repair the node that used to follow "at". */
    if (at->next != NULL) {
        at->next->prev = new_node;
    }
    at->next = new_node;

    return 1;
}
\end{lstlisting}

\begin{verbatim}
   Inserting value=5 right after the 7-node.  BEFORE and AFTER:

   BEFORE:                          AFTER:
   +------+   +------+   +------+  +------+   +------+   +------+   +------+
   | val=3|   | val=7 |   | val=9|  | val=3|   | val=7 |   | val=5|   | val=9|
   | prev |<->| prev  |<->| prev |  | prev |<->| prev  |<->| prev |<->| prev |
   | next |   | next  |   | next |  | next |   | next  |   | next |   | next |
   +------+   +------+   +------+  +------+   +------+   +------+   +------+
                  ^ at                            ^ at

   The new node (value=5) is spliced in between 7 and 9.
\end{verbatim}

Deleting a node given only its address is the doubly linked list's
signature advantage:

\begin{lstlisting}[caption={Delete a node in O(1), given its address.}]
void dlist_remove(dnode *victim) {
    /* Repair the left neighbor, if any. */
    if (victim->prev != NULL) {
        victim->prev->next = victim->next;
    }
    /* Repair the right neighbor, if any. */
    if (victim->next != NULL) {
        victim->next->prev = victim->prev;
    }
    free(victim);
}
\end{lstlisting}

No search is needed: the victim itself carries the arrows to its two
neighbors, and each neighbor points back at the victim. Two assignments
unlink it. This is why hash-table buckets, LRU caches, and memory
allocators all use doubly linked lists internally.

\subsection{The Circular Linked List}

Make the tail point back at the head instead of at \code{NULL}, and the
list becomes a ring. With a single \code{tail} pointer, both push-front
and push-back are \code{O(1)}, and a walk never ends---which is exactly
what a round-robin scheduler wants.

\begin{lstlisting}[caption={The circular node is just a node; the ring is a convention.}]
typedef struct cnode {
    int value;
    struct cnode *next;
} cnode;
\end{lstlisting}

\begin{verbatim}
   A circular list. The last node's next points back at the first.

   tail
    |
    v
   +-------+      +-------+      +-------+
   | val=1 |      | val=3 |      | val=7 |
   | next -----> | next -----> | next --+
   +-------+      +-------+      +-------+
       ^                             |
       +-----------------------------+

   - tail->next is the first node (O(1) "front").
   - tail itself is the last node (O(1) "back").
   - Walking: cur = cur->next revisits 1 after 7, forever.

   Round-robin across 3 processes:
     process 1 runs, then 2, then 3, then back to 1...
     cur = cur->next;  -- a single pointer, a single step.
\end{verbatim}

Inserting at the front of a circular list is the singly linked push
plus one extra fix-up to keep the ring closed.

\begin{lstlisting}[caption={Push at the front of a circular list.}]
int clist_push_front(cnode **tail, int value) {
    cnode *new_node;

    new_node = malloc(sizeof(cnode));
    if (new_node == NULL) {
        return 0;
    }
    new_node->value = value;

    if (*tail == NULL) {
        /* First node: the ring is just itself. */
        new_node->next = new_node;
        *tail = new_node;
        return 1;
    }

    new_node->next = (*tail)->next;   /* new node points at the old first */
    (*tail)->next = new_node;         /* tail still points at the last    */
    return 1;
}
\end{lstlisting}

\begin{verbatim}
   Why "tail" and not "head"? With a tail pointer:

     front = tail->next     O(1)
     back  = tail           O(1)

   With only a head pointer, finding the last node to push at the back
   would require walking the whole ring. The tail pointer buys O(1) on
   both ends at the cost of one field.
\end{verbatim}

\begin{insight}[The sentinel head]
The empty ring is the only case that needs care. Our push handles it
with an explicit \code{if (*tail == NULL)}. An alternative that
eliminates the branch entirely is a permanent \emph{sentinel} node that
is never data---the ring always exists, and ``empty'' means the
sentinel points at itself. Linux kernel lists (next subsection) are
built exactly this way, so that the empty list and the non-empty list
are the same code path.
\end{insight}

\subsection{The Intrusive Linked List}

Every list so far owned its data: the node \emph{was} the value. The
\emph{intrusive} list inverts this: your data structure contains the
links. Instead of \code{struct node \{ int value; node *next; \}}, you
embed a generic \code{list\_head} inside your own struct, and the list
machinery manipulates only the embedded links. The data structure never
knows what its elements are---it owns no \code{malloc}, it frees
nothing, it just arranges arrows.

\begin{lstlisting}[caption={The generic link: no data, just arrows.}]
typedef struct list_head {
    struct list_head *next;
    struct list_head *prev;
} list_head;
\end{lstlisting}

\begin{lstlisting}[caption={A struct that carries a link inside it.}]
typedef struct process {
    int pid;
    int priority;
    list_head links;     /* the intrusive hook: the list arranges this */
} process;
\end{lstlisting}

\begin{verbatim}
   Intrusive list: the chain is made of list_head fields EMBEDDED in
   the processes.

   +----------------+---------------+
   | process { pid = 1000, prio 3 } |
   |   links.next --+               |
   |   links.prev <-- NULL          |
   +----------------+---------------+
                    |   (next: down, prev: up)
                    v
   +----------------+---------------+
   | process { pid = 1001, prio 1 } |
   |   links.next --+               |
   |   links.prev <-+               |
   +----------------+---------------+
                    |
                    v
   +----------------+---------------+
   | process { pid = 1002, prio 5 } |
   |   links.next --> NULL          |
   |   links.prev <-+               |
   +----------------+---------------+

   The list code sees only the links; it never knows pid exists.
   Given a pointer to a link, the process is recovered by arithmetic:

       process *p = (process *) ((char *) link_ptr - offsetof(process, links));
\end{verbatim}

That last step, going from a member pointer back to the containing
struct, is the entire trick, and the Linux kernel gives it a name:
\code{container\_of}. It is just pointer subtraction with
\code{offsetof}.

\begin{lstlisting}[caption={container\_of: from a member back to its container.}]
#include <stddef.h>

/* Given a pointer to a member inside a struct, return the struct. */
#define container_of(ptr, type, member) \
    ((type *) ((char *) (ptr) - offsetof(type, member)))
\end{lstlisting}

The list operations are uniform and tiny. We use a permanent sentinel
\code{head} that forms a ring with itself when empty, which removes
every ``empty list'' special case.

\begin{lstlisting}[caption={Circular intrusive list with a sentinel head.}]
/* A head is a link too; it never carries data. */
void list_init(list_head *head) {
    head->next = head;
    head->prev = head;
}

/* Insert a new link between prev and next. */
static void list_add_between(list_head *new_node, list_head *prev,
                             list_head *next) {
    next->prev = new_node;
    new_node->next = next;
    new_node->prev = prev;
    prev->next = new_node;
}

/* Insert at the front of the list headed by head. */
void list_push_front(list_head *head, list_head *new_node) {
    list_add_between(new_node, head, head->next);
}

/* Unlink a link from its list. */
void list_remove(list_head *entry) {
    entry->prev->next = entry->next;
    entry->next->prev = entry->prev;
    entry->next = entry;   /* leave the link self-contained (safe) */
    entry->prev = entry;
}

/* Walk the list; on each pass you recover the container with container_of. */
#define list_for_each(pos, head) \
    for (pos = (head)->next; pos != (head); pos = pos->next)
\end{lstlisting}

\begin{verbatim}
   The sentinel head forms the ring:

   empty:  head.next = head, head.prev = head

            +-------+
            | head  |<-+
            | next--|--+
            | prev--|--+
            +-------+

   after pushing two processes: head sits between the front and the
   back, and the ring is closed through it. Walking stops when you
   return to the sentinel, not at NULL.

   Why intrusive?  Because the list is FREE: no allocation per element,
   no node ownership, and an element can belong to several lists at
   once (e.g. one "all processes" list and one "ready queue"). The cost
   is that container_of is your responsibility, and that an element
   cannot outlive its containers.
\end{verbatim}

\begin{insight}[Why Linux uses intrusive lists everywhere]
A kernel may have millions of tasks, and each task can sit in several
lists simultaneously. An intrusive list adds zero allocations and zero
per-element nodes, and the link is wherever you put it. This is the
same philosophy behind \code{container\_of}: ownership and layout
belong to the user, the list only arranges arrows.
\end{insight}

\subsection{The Unrolled Linked List}

A plain linked list scatters one small node per element across the
heap, and each \code{next} dereference may miss the cache. An
\emph{unrolled} list trades a little insertion work for locality: each
node holds a small \emph{array} of values plus one \code{next} pointer,
so a walk visits many elements per node and per cache line.

\begin{lstlisting}[caption={The unrolled node: a block plus one arrow.}]
#define BLOCK_CAP 4

typedef struct unrolled_node {
    int count;                 /* how many slots are currently live */
    int data[BLOCK_CAP];       /* the fixed-size block              */
    struct unrolled_node *next;
} unrolled_node;
\end{lstlisting}

\begin{verbatim}
   An unrolled list holding 1,2,3,4,5,6:

   +-----------------------------------------------+
   | count=4 | 1 | 2 | 3 | 4 | next ----------+    |
   +-----------------------------------------------+ |
   +-----------------------------------------------+ |
   | count=2 | 5 | 6 | _ | _ | next -> NULL     |<--+
   +-----------------------------------------------+

   - Four elements arrive in ONE node: one allocation, one cache line
     (or two), four dereferences avoided.
   - A block that fills to capacity is SPLIT into two half-full blocks.
   - Deletion can MERGE adjacent half-empty blocks to stay dense.
\end{verbatim}

Insertion walks to the right block, then either places the value in a
free slot or splits the block first.

\begin{lstlisting}[caption={Insert at a 0-based index, splitting when full.}]
int ulist_insert(unrolled_node **head, int index, int value) {
    unrolled_node *cur;
    unrolled_node *block;
    size_t pos;
    size_t i;

    if (*head == NULL) {
        block = malloc(sizeof(unrolled_node));
        if (block == NULL) {
            return 0;
        }
        block->count = 0;
        block->next = NULL;
        *head = block;
    }

    /* Advance to the block containing position "index". */
    cur = *head;
    pos = (size_t) index;
    while (cur->next != NULL && pos > (size_t) cur->count) {
        pos = pos - (size_t) cur->count;
        cur = cur->next;
    }

    /* A full block is split: move the upper half to a fresh block. */
    if (cur->count == BLOCK_CAP) {
        block = malloc(sizeof(unrolled_node));
        if (block == NULL) {
            return 0;
        }
        block->count = BLOCK_CAP / 2;
        for (i = 0; i < (size_t) block->count; i++) {
            block->data[i] = cur->data[BLOCK_CAP / 2 + i];
        }
        block->next = cur->next;
        cur->next = block;
        cur->count = BLOCK_CAP / 2;

        /* The insert belongs to the new block if it passed the middle. */
        if (pos > (size_t) cur->count) {
            pos = pos - (size_t) cur->count;
            cur = block;
        }
    }

    /* Shift the block's tail right by one, then store the value. */
    for (i = cur->count; i > pos; i--) {
        cur->data[i] = cur->data[i - 1];
    }
    cur->data[pos] = value;
    cur->count = cur->count + 1;

    return 1;
}
\end{lstlisting}

\begin{verbatim}
   Inserting 9 into the middle of a full block:

   BEFORE                                AFTER SPLIT + INSERT
   +---------------------------+         +---------------------------+
   | count=4 | 1 | 2 | 3 | 4 |          | count=2 | 1 | 2 |   |   | ---> block
   +---------------------------+         +---------------------------+
                                           count=3 | 9 | 3 | 4 |   |
                                         (new block holds the tail)

   The block no longer has to move when the list grows a bit more;
   locality is preserved at the cost of O(block) shifting per insert.
\end{verbatim}

The unrolled list shines for queues and work-stealing stacks where
traversal dominates, and it is the ancestor of the \emph{block} idea
inside B-trees (Section~4.6). It also demonstrates the Loop Counter
Rule in a shift loop where the counter must count \emph{down} without
underflowing: \code{for (i = cur->count; i > pos; i--)} stops at
\code{i == pos}, never below.

\subsection{The XOR Linked List}

A doubly linked list stores two pointers per node. The \emph{XOR
list} stores one integer: the XOR of the addresses of its two
neighbors. From one neighbor and the stored value you always recover
the other, because \code{(a \^{} b) \^{} a == b}.

\begin{lstlisting}[caption={The XOR node: one integer, two neighbors recovered.}]
#include <stdint.h>

typedef struct xor_node {
    int value;
    uintptr_t link;   /* prev XOR next, stored as an unsigned address */
} xor_node;
\end{lstlisting}

\begin{verbatim}
   head ----> [3] ----> [7] ----> [9] ----> NULL

   stored in each node:
     node(3).link = NULL  ^ node(7)     = node(7)
     node(7).link = node(3) ^ node(9)
     node(9).link = node(7) ^ NULL      = node(7)

   walking forward you know the previous node, so:
     next = node.link  XOR  prev

   walking backward you know the next node, so:
     prev = node.link  XOR  next

   It is a doubly linked list with a single pointer per node.
\end{verbatim}

Insertion at the front must repair the old head's link, and every walk
must carry the predecessor explicitly.

\begin{lstlisting}[caption={Push at the front of an XOR list.}]
xor_node *xlist_push_front(xor_node *head, int value) {
    xor_node *new_node;

    new_node = malloc(sizeof(xor_node));
    if (new_node == NULL) {
        return NULL;
    }
    new_node->value = value;

    /* new_node->link must be: prev(NULL) XOR next(head) = head. */
    new_node->link = (uintptr_t) head;

    /* The old head's link was (0 XOR its_next); it becomes
       (its_next XOR new_node), so that its_next is still recoverable. */
    if (head != NULL) {
        head->link = head->link ^ (uintptr_t) new_node;
    }

    return new_node;
}
\end{lstlisting}

\begin{lstlisting}[caption={Walk forward, remembering the predecessor.}]
void xlist_print(const xor_node *head) {
    const xor_node *cur;
    const xor_node *prev;
    const xor_node *next;
    uintptr_t stored;

    prev = NULL;
    cur = head;
    while (cur != NULL) {
        printf("%d ", cur->value);

        stored = cur->link;
        next = (const xor_node *) (stored ^ (uintptr_t) prev);

        prev = cur;
        cur = next;
    }
    printf("\n");
}
\end{lstlisting}

\begin{verbatim}
   The walk keeps TWO pointers alive: the current node and the one
   before it. From those two, the next node is arithmetic, not a
   dereference:

     state:   prev        cur         (stored link)      next
              NULL        [3]         node(7)     -->    [7]
              [3]         [7]         3^9         -->    [9]
              [7]         [9]         node(7)     -->    NULL
\end{verbatim}

\begin{ruleofthumb}[XOR lists are a curiosity, use them sparingly]
The XOR list saves one pointer per node, and pays for it in three ways:
the walk always needs the neighbor (you cannot jump into the middle of
a list and walk from there); casting pointers to \code{uintptr\_t} and
back is implementation-defined, so the trick is not portable to every
platform; and it is incompatible with garbage collectors and with any
form of bounds checking that inspects pointer provenance. Admire it,
understand the arithmetic, then prefer a conventional doubly linked
list in production code.
\end{ruleofthumb}

\section{The Skip List}

The skip list is the structure that made Salvatore Sanfilippo's
Redis famous as a sorted-set implementation, and it deserves its
reputation. It is a stack of linked lists: the bottom level holds every
element, and each higher level holds a random subset, spaced widely
enough that searches leap over large ranges and land near their target
in \code{O(log n)} expected time.

\begin{verbatim}
   A skip list over 3,9,12,17,21,25. Towers run vertically.

   level 4:  head ----------------------------------------> NULL
   level 3:  head ----------------------> 17 ------------> NULL
   level 2:  head ---------> 9 --------> 17 ------------> NULL
   level 1:  head --> 3 ---> 9 --------> 17 ---> 21 ---> NULL
   level 0:  head --> 3 --> 9 --> 12 --> 17 --> 21 --> 25 -> NULL

   Searching for 21, starting at the head's highest level:
     level 4:  next is NULL,            drop down
     level 3:  17 < 21  move; NULL,     drop down
     level 2:  17 < 21  move; NULL,     drop down
     level 1:  21 found in 2 hops here; overall ~4 hops instead of 6.

   The expected search length is O(log n): each level cuts the space
   by a constant factor, exactly like the levels of a balanced tree--
   but a skip list is a SORTED SET built from linked lists, with no
   rotations and no rebalancing passes.
\end{verbatim}

A node carries an array of forward pointers; its height is chosen
randomly, so tall nodes are rare. We build the node with the flexible
array member from Chapter~3.

\begin{lstlisting}[caption={The skip list structures.}]
#include <stdio.h>
#include <stdlib.h>
#include <string.h>

#define SKIP_MAX_LEVEL 16
#define SKIP_P          4   /* 1/P probability of one more level */

typedef struct skiplist_node {
    int value;
    int level;                        /* how many forward links this node has */
    struct skiplist_node *forward[];  /* forward[0..level-1]                 */
} skiplist_node;

typedef struct skiplist {
    skiplist_node *head;   /* tower of SKIP_MAX_LEVEL forward links */
    int level;             /* highest level currently in use        */
    int size;              /* number of elements                    */
} skiplist;

/* Random tower height: geometric with expectation 1/(1-1/P). */
static int random_level(void) {
    int level;

    level = 1;
    while ((rand() % SKIP_P) == 0 && level < SKIP_MAX_LEVEL) {
        level = level + 1;
    }
    return level;
}

void skiplist_init(skiplist *sl) {
    int i;

    sl->head = malloc(sizeof(skiplist_node) +
                      (size_t) SKIP_MAX_LEVEL * sizeof(skiplist_node *));
    if (sl->head == NULL) {
        exit(1);
    }
    for (i = 0; i < SKIP_MAX_LEVEL; i++) {
        sl->head->forward[i] = NULL;
    }
    sl->head->level = SKIP_MAX_LEVEL;
    sl->level = 1;
    sl->size = 0;
}
\end{lstlisting}

Search first, because insert and delete reuse its shape: descend level
by level, each time advancing as far right as possible without
overshooting.

\begin{lstlisting}[caption={Search: leap down and to the right.}]
/* Returns 1 if value is present. */
int skiplist_search(const skiplist *sl, int value) {
    const skiplist_node *cur;
    size_t i;

    cur = sl->head;
    for (i = sl->level; i > 0; i--) {
        while (cur->forward[i - 1] != NULL &&
               cur->forward[i - 1]->value < value) {
            cur = cur->forward[i - 1];
        }
    }

    /* After the top-down descent, the answer is the next node at level 0. */
    cur = cur->forward[0];
    if (cur != NULL && cur->value == value) {
        return 1;
    }
    return 0;
}
\end{lstlisting}

Notice the down-counting loop written safely: the \code{size\_t} counter
\code{i} counts \emph{down to 1}, never to 0, and the actual level
index is \code{i - 1}. The loop counter never underflows.

Insertion records, for each level, the last node visited before the
insertion point; those nodes' forward links are exactly where the new
tower must be spliced.

\begin{lstlisting}[caption={Insert: record the update path, then splice the tower.}]
/* Insert value. Returns 1 on success, 0 on allocation failure. */
int skiplist_insert(skiplist *sl, int value) {
    skiplist_node *update[SKIP_MAX_LEVEL];
    skiplist_node *cur;
    skiplist_node *new_node;
    size_t i;
    int new_level;

    /* Descend, remembering the last node at each level. */
    cur = sl->head;
    for (i = sl->level; i > 0; i--) {
        while (cur->forward[i - 1] != NULL &&
               cur->forward[i - 1]->value < value) {
            cur = cur->forward[i - 1];
        }
        update[i - 1] = cur;
    }

    /* Choose the tower height; extend the head if we got taller. */
    new_level = random_level();
    if (new_level > sl->level) {
        for (i = sl->level; i < (size_t) new_level; i++) {
            update[i] = sl->head;
        }
        sl->level = new_level;
    }

    new_node = malloc(sizeof(skiplist_node) +
                      (size_t) new_level * sizeof(skiplist_node *));
    if (new_node == NULL) {
        return 0;
    }
    new_node->value = value;
    new_node->level = new_level;

    /* Splice the new node into each level of the update path. */
    for (i = 0; i < (size_t) new_level; i++) {
        new_node->forward[i] = update[i]->forward[i];
        update[i]->forward[i] = new_node;
    }

    sl->size = sl->size + 1;
    return 1;
}
\end{lstlisting}

\begin{verbatim}
   Inserting 10 into the list drawn above, with a 2-tall tower:

   level 2:  update[2] = head      level 2:  head -------------------> 17 -> ...
                                             ^
                                             |new
   level 1:  update[1] = 9         level 1:  head --> 9 ---> [10] ---> 17 -> ...
   level 0:  update[0] = 9         level 0:  head -> 3 -> 9 -> [10] -> 12 -> ...

   update[k] is the node whose forward[k] must point at the new node.
   After splicing, 10 is reachable at every level of its own tower.
\end{verbatim}

Deletion is the mirror image: record the update path, and for each
level of the victim's tower, route \code{update[i]->forward[i]} around
the victim; then free it.

\begin{verbatim}
   The skip list completes the "all the lists" tour:

   structure       search     insert   delete    memory/node
   singly          O(n)       O(1)+    O(n)      1 pointer
   doubly          O(n)       O(1)+    O(1)+     2 pointers
   circular        O(n)       O(1)+    O(1)+     ring
   intrusive       O(n)       O(1)+    O(1)+     embedded, free
   unrolled        O(n)       O(n)     O(n)      block of values
   XOR             O(n)       O(1)+    O(1)+     1 pointer, needs neighbor
   skip list       O(log n)*  O(log n)* O(log n)* towers (probabilistic)
   * expected, and with the same code for sorted ranges.
\end{verbatim}

\begin{insight}[Lists beat arrays for one thing: stable pointers]
An array moves its elements on insert and delete; a list does not. The
skip list keeps that property and adds logarithmic search. That is why
Redis implements sorted sets with skip lists rather than balanced
trees: the code is dramatically simpler, the pointers stay stable, and
the expected performance matches. When you must keep addresses valid
while sorting and searching, reach for a skip list.
\end{insight}

\begin{researchnote}[The skip list's origin]
The skip list was invented by William Pugh in 1990, and its original
paper is worth reading because it makes the central claim of this
section explicit: a sorted structure built from linked lists can match
a balanced tree's guarantees without the tree's machinery.

\begin{verbatim}
   An extract from the paper's abstract:

       "Skip lists are a probabilistic data structure that seem likely
        to supplant balanced trees as the implementation method of
        choice for many applications, since skip lists are
        substantially simpler to implement than balanced trees, and
        provide expected search time proportional to log n."

   In plain terms: "probabilistic" means the towers are chosen by
   coin flips rather than by enforcing a strict balance invariant;
   "expected search time proportional to log n" is the same bound a
   balanced tree guarantees, but the skip list achieves it with far
   less bookkeeping.
\end{verbatim}

The practical payoff is visible in this book's own example: Redis, a
production database serving billions of operations, uses exactly this
structure for sorted sets. A data structure invented in a research
paper is not a classroom curiosity; it is running in production all
around you.

\paperref{W.~Pugh, ``Skip Lists: A Probabilistic Alternative to Balanced
Trees,'' Communications of the ACM 33(6), 1990}{https://15721.courses.cs.cmu.edu/spring2018/papers/08-oltpindexes1/pugh-skiplists-cacm1990.pdf}
\end{researchnote}

\section{Hash Tables}

A hash table trades a function for a search. Instead of comparing a key
against every stored key, we compute a number from the key---the
\emph{hash}---reduce it modulo the table size to get a \emph{bucket
index}, and go straight there. Perfect hashing would give
\code{O(1)}. Reality gives \code{O(1)} expected, provided we handle
collisions (two different keys hashing to the same bucket) and keep the
table from filling up.

\begin{verbatim}
   The core idea:

      key "abc" --hash--> 281872  --mod 8-->  bucket 0
      key "xyz" --hash--> 450817  --mod 8-->  bucket 1
      key "ghi" --hash--> 226514  --mod 8-->  bucket 2

      buckets
        +------+------+------+------+------+------+------+------+
        | [0]  | [1]  | [2]  | [3]  | [4]  | [5]  | [6]  | [7]  |
        +------+------+------+------+------+------+------+------+
         key     key    key
         "abc"   "xyz"  "ghi"

   Two keys with the same bucket collide; the two major collision
   strategies are the two subsections below.
\end{verbatim}

The hash function must be fast, must spread keys evenly, and must be
\emph{deterministic}. The classic djb2 reads a string and mixes it with
\code{h = h * 33 + c}.

\begin{lstlisting}[caption={djb2: a simple, proven string hash.}]
static unsigned int hash_str(const char *s) {
    unsigned int h;
    int c;

    h = 5381;
    c = (unsigned char) *s;
    while (c != 0) {
        h = (h * 33) + (unsigned int) c;
        s = s + 1;
        c = (unsigned char) *s;
    }
    return h;
}
\end{lstlisting}

Note the cast of each character to \code{unsigned char}: \code{char}
may be signed, and negative bytes would otherwise corrupt the mix.
The multiplication by 33 and the accumulation spread the bits across
the whole word, which is what turns similar keys into dissimilar
bucket indexes.

\subsection{Separate Chaining}

Separate chaining makes each bucket a small linked list. A collision is
not an error; it just makes the bucket's chain one node longer. The
whole table is an array of chain heads.

\begin{lstlisting}[caption={A chained hash table.}]
#include <stdio.h>
#include <stdlib.h>
#include <string.h>

typedef struct ht_node {
    char *key;                /* owned copy of the key */
    int   value;
    struct ht_node *next;     /* next node in this bucket */
} ht_node;

typedef struct hashtable {
    ht_node **buckets;        /* array of chain heads */
    int nbuckets;             /* current number of buckets */
    int used;                 /* number of live entries   */
} hashtable;
\end{lstlisting}

\begin{verbatim}
   buckets is an array of pointers; each pointer is the head of a chain.

        +------+      +------+      +------+      +------+
        | [0]  |      | [1]  |      | [2]  |      | [3]  |
        +------+      +------+      +------+      +------+
          |              |             |             |
          v              v             v             v
        +-------+      +-------+     +-------+     +-------+
        | "abc" |      | "ghi" |     | "def" |     | "xyz" |
        | 42    |      | 7     |     | 99    |     | 13    |
        | next -->NULL | next -->NULL | next -->NULL | next -->NULL
        +-------+      +-------+     +-------+     +-------+
        +-------+
        | "zzz" |      <-- "zzz" collided with "ghi": the chain grows
        | 1     |
        | next -->NULL
        +-------+

   A collision simply lengthens a chain. Load factor = used/nbuckets;
   chains average exactly the load factor.
\end{verbatim}

Lookup walks the bucket's chain comparing with \code{strcmp}; update
replaces a value when the key is found; otherwise a new node is pushed
onto the chain.

\begin{lstlisting}[caption={Insert or update. Returns 1 on success.}]
int ht_set(hashtable *ht, const char *key, int value) {
    unsigned int idx;
    ht_node *cur;
    ht_node *new_node;

    idx = hash_str(key) % (unsigned int) ht->nbuckets;

    /* Update in place if the key already exists in its bucket. */
    cur = ht->buckets[idx];
    while (cur != NULL) {
        if (strcmp(cur->key, key) == 0) {
            cur->value = value;
            return 1;
        }
        cur = cur->next;
    }

    /* Not present: allocate a node and its owned key copy. */
    new_node = malloc(sizeof(ht_node));
    if (new_node == NULL) {
        return 0;
    }
    new_node->key = strdup(key);
    if (new_node->key == NULL) {
        free(new_node);
        return 0;
    }
    new_node->value = value;

    new_node->next = ht->buckets[idx];
    ht->buckets[idx] = new_node;

    ht->used = ht->used + 1;
    return 1;
}
\end{lstlisting}

\begin{lstlisting}[caption={Lookup. Returns 1 and sets *out when found.}]
int ht_get(const hashtable *ht, const char *key, int *out) {
    unsigned int idx;
    const ht_node *cur;

    idx = hash_str(key) % (unsigned int) ht->nbuckets;

    cur = ht->buckets[idx];
    while (cur != NULL) {
        if (strcmp(cur->key, key) == 0) {
            *out = cur->value;
            return 1;
        }
        cur = cur->next;
    }
    return 0;
}
\end{lstlisting}

The killer detail in \code{ht\_set}: the key is \emph{copied} with
\code{strdup}, so the table owns the key for its whole life and the
caller may reuse or free the original buffer. Ownership is explicit on
both sides. The destroy function then frees every node and every owned
key copy.

When the load factor climbs (say past 0.75), we \emph{grow} the table:
allocate a new, larger bucket array and rehash every live node into it.
Rehashing matters---\code{idx} depends on \code{nbuckets}, so the old
bucket number is meaningless in the new array.

\begin{lstlisting}[caption={Grow the table and rehash every node.}]
static int ht_resize(hashtable *ht, int new_count) {
    ht_node **old_buckets;
    ht_node **new_buckets;
    ht_node *cur;
    ht_node *next;
    unsigned int idx;
    size_t i;

    new_buckets = calloc((size_t) new_count, sizeof(ht_node *));
    if (new_buckets == NULL) {
        return 0;
    }

    old_buckets = ht->buckets;
    for (i = 0; i < (size_t) ht->nbuckets; i++) {
        cur = old_buckets[i];
        while (cur != NULL) {
            next = cur->next;                /* save before relinking */

            idx = hash_str(cur->key) % (unsigned int) new_count;
            cur->next = new_buckets[idx];    /* rehang into the new table */
            new_buckets[idx] = cur;

            cur = next;
        }
    }

    free(old_buckets);
    ht->buckets = new_buckets;
    ht->nbuckets = new_count;
    return 1;
}
\end{lstlisting}

Because the nodes themselves are reused---only their \code{next}
pointers are rewired---the resize is a cheap pointer shuffle, not a
reallocation of every key.

\begin{ruleofthumb}[Keep the load factor bounded]
Chaining degrades gracefully, but long chains defeat the \code{O(1)}
promise. Grow when \code{used / nbuckets} exceeds roughly 0.75. The
resize doubles the array, so the amortized cost of growth is a
constant per insert---the same argument that pays for dynamic arrays.
\end{ruleofthumb}

\subsection{Open Addressing}

Open addressing keeps all entries \emph{inside} the array itself:
no chains, no per-node allocations, better cache locality. A collision
is resolved by \emph{probing} forward to the next free slot. The
simplest probe is \emph{linear}: try slot \code{i}, then \code{i+1},
then \code{i+2}, wrapping around the end of the array.

\begin{lstlisting}[caption={An open-addressing table with linear probing.}]
typedef enum slot_state {
    SLOT_EMPTY,       /* never used: the search can stop here */
    SLOT_USED,        /* holds a live key */
    SLOT_TOMBSTONE    /* was used, then deleted: search continues */
} slot_state;

typedef struct oa_table {
    char       **keys;
    int        *values;
    slot_state *states;
    int        capacity;
    int        used;
} oa_table;
\end{lstlisting}

\begin{verbatim}
   capacity = 8.  hash("abc")%8 = 3, hash("xyz")%8 = 4, hash("ghi")%8 = 5,
   hash("def")%8 = 3  (collides with "abc").

   Insert in order: "abc", "xyz", "ghi", "def".

   slot:  0        1        2        3        4        5        6        7
        +--------+--------+--------+--------+--------+--------+--------+--------+
        | EMPTY  | EMPTY  | EMPTY  | "abc"  | "xyz"  | "ghi"  | "def"  | EMPTY  |
        |        |        |        |  42    |   7    |   1    |  99    |        |
        +--------+--------+--------+--------+--------+--------+--------+--------+

   "def" wanted slot 3, found it full, probed 4, full, 5, full, 6, empty.
   Linear probing turns collisions into a short run of nearby probes.
   Deletion leaves a TOMBSTONE so later searches keep walking past the
   hole instead of stopping at a false "not found".
\end{verbatim}

Insertion probes forward, stopping at the first empty or tombstone
slot; if it finds the key already present, it updates instead.

\begin{lstlisting}[caption={Insert with linear probing. Returns 1 on success.}]
int oa_set(oa_table *t, const char *key, int value) {
    unsigned int idx;
    size_t i;
    int slot;

    idx = hash_str(key) % (unsigned int) t->capacity;

    for (i = 0; i < (size_t) t->capacity; i++) {
        slot = (int) ((idx + (unsigned int) i) % (unsigned int) t->capacity);

        if (t->states[slot] == SLOT_EMPTY ||
            t->states[slot] == SLOT_TOMBSTONE) {
            /* Free slot: the probe stops here. */
            if (t->keys[slot] != NULL) {
                free(t->keys[slot]);
            }
            t->keys[slot] = strdup(key);
            if (t->keys[slot] == NULL) {
                return 0;
            }
            t->values[slot] = value;
            t->states[slot] = SLOT_USED;
            t->used = t->used + 1;
            return 1;
        }

        if (t->states[slot] == SLOT_USED && strcmp(t->keys[slot], key) == 0) {
            t->values[slot] = value;      /* update an existing key */
            return 1;
        }
    }

    return 0;   /* table full */
}
\end{lstlisting}

Lookup must probe exactly the same way, and it must not stop at a
tombstone---the key it seeks may be past it. It stops only at an empty
slot, which is the proof the key is absent.

\begin{lstlisting}[caption={Lookup with the same probe sequence.}]
int oa_get(const oa_table *t, const char *key, int *out) {
    unsigned int idx;
    size_t i;
    int slot;

    idx = hash_str(key) % (unsigned int) t->capacity;

    for (i = 0; i < (size_t) t->capacity; i++) {
        slot = (int) ((idx + (unsigned int) i) % (unsigned int) t->capacity);

        if (t->states[slot] == SLOT_EMPTY) {
            return 0;        /* the key is definitely not present */
        }
        if (t->states[slot] == SLOT_USED && strcmp(t->keys[slot], key) == 0) {
            *out = t->values[slot];
            return 1;
        }
    }

    return 0;
}
\end{lstlisting}

\begin{verbatim}
   Looking up "def" (hash %8 = 3):

     slot 3: USED "abc"  -- not a match, keep probing
     slot 4: USED "xyz"  -- not a match
     slot 5: USED "ghi"  -- not a match
     slot 6: USED "def"  -- match! return 99

   Looking up "qrs" (hash %8 = 3):

     slot 3..5: used, no match
     slot 6:    USED "def", no match
     slot 7:    EMPTY  -- stop. "qrs" is absent. We cannot skip slot 7;
     if we stopped at the first non-EMPTY we would report "absent" for
     keys that live further right. That is why tombstones exist.
\end{verbatim}

\begin{ruleofthumb}[Open addressing must stay loose]
Probing chains merge into clusters as the table fills, and long
clusters make every lookup slower. Keep the load factor below about
0.7, grow by reinserting all entries into a doubled array, and clear
tombstones on resize so they stop accumulating. Clustering is the
reason variants such as quadratic probing and double hashing exist:
they spread probes so clusters do not form.
\end{ruleofthumb}

\begin{verbatim}
   chaining vs open addressing:

   concern            separate chaining        open addressing
   ----------------   --------------------     ----------------------
   memory             nodes + bucket array     one flat array
   allocation         one per insert           none during insert
   locality           poor (chains scatter)    excellent (same array)
   deletion           free the node            tombstone + cleanup
   safe load factor   up to ~1.0               must stay < 0.7
   worst case         long chain               cluster / full table
   best for           generic, unknown keys    hot, cache-bound tables
\end{verbatim}

\begin{researchnote}[Collision handling has been studied since 1956]
The choice between collision strategies---our chaining versus our
linear probing---is not folklore. It was studied quantitatively in one
of the very first papers on hash tables, by E.~J.~Isaac and
R.~C.~Singleton, published in 1956, and the core trade-off they
identified is still the one you just saw in code.

\begin{verbatim}
   What the 1956 study established (paraphrased, not quoted):

       "The number of probes required to find an item depends on how
        the table is filled and on the collision technique used;
        comparing the techniques requires computing this cost."

   In our terms: chaining keeps clusters short by making the bucket a
   chain (cost grows with load factor), while open addressing pays a
   probe penalty that worsens as clusters form. That distinction is the
   load-factor table above, and it has not changed in seventy years.
\end{verbatim}

Why cite a paper from 1956? Because it shows that the design questions
you are making choices about today---probe sequences, load factors,
memory vs.\ locality---are the same questions the field has measured
since its beginning. The details have evolved (double hashing, Robin
Hood hashing), but the fundamental trade-off is a research result, not
a guess.

\paperref{E.~J.~Isaac, R.~C.~Singleton, ``An Analysis of Alternative
Collision Detection Techniques,'' Communications of the ACM, 1956}{https://dl.acm.org/doi/10.1145/320604.320609}
\end{researchnote}

\section{Trees}

\subsection{The General (N-ary) Tree}

A tree is a set of nodes, each with a value and a set of \emph{children}
that are themselves trees. When a node may have any number of children,
the simplest encoding carries two pointers: one to the first child, one
to the next sibling. The child list is a linked list, so the tree is a
list of lists of trees.

\begin{lstlisting}[caption={The N-ary node: first child, next sibling.}]
typedef struct tree_node {
    int value;
    struct tree_node *first_child;   /* head of the children list */
    struct tree_node *next_sibling;  /* next sibling in that list */
} tree_node;
\end{lstlisting}

\begin{verbatim}
   The tree:

                    (1)
                   / | \
                 (2)(3)(4)
                 / \    |
               (5) (6) (7)

   Encoded with first_child / next_sibling:

   (1) -- first_child --> (2) -- next_sibling --> (3) -- next_sibling --> (4)
                              |                      |                    |
                              v                      v                    v
                      +---------+               (NULL)                  (7)
                      | (5)     |
                      | next_sib --+            (3) has no children:
                      +------------+            first_child = NULL
                                   v
                              +---------+
                              | (6)     |
                              | next_sib --> NULL
                              +---------+

   Reading the picture: every node's children hang under it in a
   singly linked chain; the chain's links are "next_sibling".
\end{verbatim}

A preorder traversal visits a node, then each of its children in turn;
because the children form a list, the traversal is a loop over siblings
recursing into each child.

\begin{lstlisting}[caption={Preorder traversal of an N-ary tree.}]
void tree_print_preorder(const tree_node *root) {
    const tree_node *child;

    if (root == NULL) {
        return;
    }

    printf("%d ", root->value);

    child = root->first_child;
    while (child != NULL) {
        tree_print_preorder(child);
        child = child->next_sibling;
    }
}
\end{lstlisting}

This encoding is how compilers hold syntax trees (an operator with an
arbitrary number of operands), how file systems hold directory
hierarchies, and how HTML/XML documents are represented in memory.
Its cost is that a node needs a \emph{list} to reach its children; its
gift is that children of one node never share storage with anyone.

\subsection{Binary Trees and Binary Search Trees}

A binary tree is the special case where every node has at most two
children, conventionally named \code{left} and \code{right}. It is the
most studied shape in computing, because a \emph{binary search tree}
(BST) with the ordering property---left subtree holds smaller values,
right subtree holds larger---gives \code{O(h)} search, where
\code{h} is the tree height.

\begin{lstlisting}[caption={The BST node.}]
typedef struct bst_node {
    int value;
    struct bst_node *left;
    struct bst_node *right;
} bst_node;
\end{lstlisting}

\begin{verbatim}
   Inserting 3,1,4,2,5 in that order:

                3
               / \
              1   4
               \   \
                2   5

   The ordering property, stated at every node:
     everything in the left subtree  is smaller than the node;
     everything in the right subtree is larger than the node.
   Because of it, search is a descent: at each node, go left, go right,
   or stop. The value 2 is found in three steps, not five.
\end{verbatim}

Insertion reuses the pointer-to-pointer technique from the linked
lists: walk a \code{bst\_node **} down until it names an empty link,
then hang the new leaf there.

\begin{lstlisting}[caption={Insert: descend to an empty link and hang a leaf.}]
int bst_insert(bst_node **root, int value) {
    bst_node *new_node;
    bst_node **cur;

    cur = root;
    while (*cur != NULL) {
        if (value < (*cur)->value) {
            cur = &(*cur)->left;
        } else if (value > (*cur)->value) {
            cur = &(*cur)->right;
        } else {
            return 1;       /* already present: nothing to do */
        }
    }

    new_node = malloc(sizeof(bst_node));
    if (new_node == NULL) {
        return 0;
    }
    new_node->value = value;
    new_node->left = NULL;
    new_node->right = NULL;

    *cur = new_node;        /* the empty link now points at the leaf */
    return 1;
}
\end{lstlisting}

\begin{verbatim}
   Inserting 2 into the tree above. The descent:

     cur = &root        (names the 3-node)
     2 < 3  -> cur = &(3.left)     (names the 1-node)
     2 > 1  -> cur = &(1.right)    (names NULL -- an empty link)

   *cur = new_node hangs the leaf; no parent variable, no special case.

                3
               / \
              1   4
             / \   \
            ... 2   5
\end{verbatim}

Search is the same descent, returning true when the value is found.

\begin{lstlisting}[caption={Search: the descent in miniature.}]
int bst_search(const bst_node *root, int value) {
    const bst_node *cur;

    cur = root;
    while (cur != NULL) {
        if (value < cur->value) {
            cur = cur->left;
        } else if (value > cur->value) {
            cur = cur->right;
        } else {
            return 1;
        }
    }
    return 0;
}
\end{lstlisting}

A BST walked in-order---left, node, right---visits every value in
sorted order, which is exactly what a sorted ``print the whole
collection'' needs.

\begin{lstlisting}[caption={In-order traversal yields sorted output.}]
void bst_print_inorder(const bst_node *root) {
    if (root == NULL) {
        return;
    }
    bst_print_inorder(root->left);
    printf("%d ", root->value);
    bst_print_inorder(root->right);
}
\end{lstlisting}

Deletion is the one operation with three cases, and it is where BSTs
earn their reputation. The three cases are a leaf, a node with one
child, and a node with two children.

\begin{lstlisting}[caption={Remove, handling the three structural cases.}]
/* Remove value. Returns 1 if it was present. */
int bst_remove(bst_node **root, int value) {
    bst_node *victim;
    bst_node *successor;

    if (*root == NULL) {
        return 0;
    }

    if (value < (*root)->value) {
        return bst_remove(&(*root)->left, value);
    }
    if (value > (*root)->value) {
        return bst_remove(&(*root)->right, value);
    }

    /* Found it: *root is the node to remove. */
    victim = *root;

    if (victim->left == NULL && victim->right == NULL) {
        /* Case 1: a leaf. Drop it. */
        free(victim);
        *root = NULL;
        return 1;
    }

    if (victim->left == NULL) {
        /* Case 2: only a right child; it replaces the node. */
        *root = victim->right;
        free(victim);
        return 1;
    }
    if (victim->right == NULL) {
        /* Case 2: only a left child; it replaces the node. */
        *root = victim->left;
        free(victim);
        return 1;
    }

    /* Case 3: two children. Steal the value of the in-order successor
       (the leftmost node of the right subtree), then delete that
       successor node, which has at most one child.               */
    successor = victim->right;
    while (successor->left != NULL) {
        successor = successor->left;
    }
    victim->value = successor->value;
    return bst_remove(&victim->right, successor->value);
}
\end{lstlisting}

\begin{verbatim}
   Case 1 (leaf):   delete 9               6
                                             / \
        6            6                    3   9      ->  3's right = empty
       / \          / \                  / \   \
      3   9   ==>  3   (empty)          1   5   (gone)
     / \          / \
    1   5        1   5

   Case 2 (one child):   delete 3         6
                                             / \
        6            6                    1   9     ->  3's right promoted
       / \          / \                  / \   \
      3   9   ==>  1   9                1   5   7
     / \   \      / \   \
    1   5   7    1   5   7

   Case 3 (two children):   delete 6
        6                   7
       / \                 / \
      3   9       ==>     3   9
     / \   \             / \   \
    1   5   7           1   5   (the 7-node is gone)
\end{verbatim}

\begin{insight}[The height is the whole story]
A BST is only fast if it stays shallow. Inserting already-sorted
values yields a right-only chain of height \code{n}, and every
operation becomes \code{O(n)}---a linked list wearing a tree costume.
Balanced variants (AVL, red--black) guarantee \code{O(log n)} by
rotating on insertion. We do not implement them here because the
B-trees below give the same guarantee with a structure the disk and
cache actually like.
\end{insight}

\subsection{B-Trees}

A B-tree is the tree that databases and file systems actually use. It
is a balanced search tree with \emph{wide} nodes: each node holds many
keys and many children, so a tree of a million keys is only a few
levels tall, and each visit to a node can read an entire block in one
I/O. The balance is absolute---every leaf sits at the same depth---and
it is maintained by splitting nodes that overflow, not by rotations.

\begin{lstlisting}[caption={A B-tree node of order 4 (max 3 keys, 4 children).}]
#define BTREE_ORDER 4

typedef struct btree_node {
    int nkeys;                     /* keys currently stored */
    int keys[BTREE_ORDER - 1];     /* up to 3 keys          */
    struct btree_node *child[BTREE_ORDER];  /* up to 4 children */
    int leaf;                      /* 1 if this node has no children */
} btree_node;
\end{lstlisting}

\begin{verbatim}
   A B-tree of order 4 ("2-3-4 tree").  Keys in [ ] , children below.

                    [  20  |  40  ]
                   /       |       \
          [ 5 | 10 ]  [ 25 | 30 | 35 ]   [ 60 | 70 ]
          /    |    \    ...              ...
       (1,3)  (7,8) (11,15)

   Invariants (order m = 4):
     - a node with k keys has k+1 children;
     - every node (except the root) holds between m/2 - 1 and m - 1 keys;
     - all leaves sit at the same depth;
     - keys inside a node are ordered, and children bracket them:
         child[i]   holds keys  <  keys[i]
         child[i+1] holds keys  >  keys[i]

   Search is one node visit per level, and a node visit is a linear
   scan of a handful of keys. With wide nodes, height stays tiny:
   order 200 and a 2-level tree already covers 40,000 keys.
\end{verbatim}

Search first, because insert descends the same path.

\begin{lstlisting}[caption={Search a B-tree, one node visit per level.}]
/* Returns 1 if key is present. */
int btree_search(const btree_node *root, int key) {
    const btree_node *cur;
    size_t i;

    cur = root;
    while (cur != NULL) {
        /* Find the first key >= key in this node. */
        i = 0;
        while (i < (size_t) cur->nkeys && key > cur->keys[i]) {
            i = i + 1;
        }

        if (i < (size_t) cur->nkeys && key == cur->keys[i]) {
            return 1;              /* found in this node */
        }
        if (cur->leaf) {
            return 0;              /* reached a leaf without finding it */
        }
        cur = cur->child[i];       /* descend into the correct child */
    }
    return 0;
}
\end{lstlisting}

Insertion never grows a node beyond \code{m - 1} keys, because we
\emph{split} any node that is full on the way down. The middle key
moves up into the parent, and the two halves become children of the
parent. This is the single idea behind all B-tree variants.

\begin{lstlisting}[caption={Split a full child of cur at child index i.}]
/* Precondition: cur->child[i] is full (holds BTREE_ORDER-1 keys). */
void btree_split_child(btree_node *cur, size_t i) {
    btree_node *child;
    btree_node *right;
    size_t mid;
    size_t j;

    child = cur->child[i];
    right = malloc(sizeof(btree_node));
    if (right == NULL) {
        exit(1);
    }

    mid = (BTREE_ORDER - 1) / 2;
    right->leaf = child->leaf;
    right->nkeys = child->nkeys - (int) mid - 1;

    /* Move the upper half of child's keys into right. */
    for (j = 0; j < (size_t) right->nkeys; j++) {
        right->keys[j] = child->keys[mid + 1 + j];
    }
    /* Move the upper half of child's children into right. */
    if (child->leaf == 0) {
        for (j = 0; j <= (size_t) right->nkeys; j++) {
            right->child[j] = child->child[mid + 1 + j];
        }
    }

    /* Make room in the parent, promote the middle key, link right in. */
    for (j = cur->nkeys; j > i; j--) {
        cur->keys[j] = cur->keys[j - 1];
        cur->child[j + 1] = cur->child[j];
    }
    cur->keys[i] = child->keys[mid];
    cur->child[i + 1] = right;

    child->nkeys = (int) mid;
    cur->nkeys = cur->nkeys + 1;
}
\end{lstlisting}

\begin{verbatim}
   Splitting a full child during insert.

   BEFORE (child at index 1 is full):          AFTER split:
        [ 20 ]                                     [ 20 | 30 ]
        /      \                                 /    |      \
   [5|10]  [25|30|35]                     [5|10]  [25]  [35]
   (full: 3 keys)

   - the middle key 30 is PROMOTED into the parent;
   - the child splits into two nodes of 1 key each (for order 4,
     a node is at minimum half full after a split);
   - the parent grew by one key, and may itself need a split next.
\end{verbatim}

Insert descends from the root, splitting every full node it meets, so
that when the new key reaches a leaf there is always room for it. The
root is the only node that may ever be split alone, and splitting the
root is the only way a B-tree gains a level---balance is never
temporarily broken.

\begin{verbatim}
   The growth of a B-tree is the height growing UPWARD:

   [1|2|3]  insert 4  ->  [1]  [2]  [3|4]   (2 moves up, root splits)
                              |
                          [ 2 ]
                         /     \
                      [1]     [3|4]   ->   insert more... height grows at
                                          the root, never in the middle.
\end{verbatim}

\begin{ruleofthumb}[B-trees trade memory for I/O]
Each node is sized to a block (a disk sector or a cache line), so one
node read fetches many keys at once. Search becomes ``read a block,
scan a few keys, repeat,'' with a path length that is the logarithm of
the key count in base \emph{order}. That is why database indexes and
filesystem metadata use B-trees and not BSTs.
\end{ruleofthumb}

\subsection{B+ Trees}

The B+ tree takes the B-tree and separates two jobs: \emph{internal
nodes hold only keys, for routing}, and \emph{all data lives in the
leaves}, which are linked together in sorted order. A lookup or range
scan descends once and then walks the leaf chain, which is sequential
reads---exactly what rotating disks and SSDs love.

\begin{lstlisting}[caption={The B+ leaf adds a sibling pointer.}]
typedef struct bplus_node {
    int leaf;                 /* 1 for a leaf, 0 for an internal node */
    int nkeys;
    int keys[BTREE_ORDER - 1];
    struct bplus_node *child[BTREE_ORDER];
    struct bplus_node *next;  /* used ONLY by leaves: the sibling link */
} bplus_node;
\end{lstlisting}

\begin{verbatim}
   A B+ tree. Internal nodes hold routing keys only; the data rows
   live in the leaves, and the leaves are a sorted linked list.

            [ 20 | 40 ]              <- internal: pure routing
           /     |      \
    [ 5|10 ]  [ 25|30 ]  [ 60|70 ]   <- internal: still routing
       |          |          |
       v          v          v
   +--------------------------------------------------------------+
   | leaf (1,3,5,7) -> leaf (20,25,30) -> leaf (40,60,70) -> NULL |
   +--------------------------------------------------------------+

   "SELECT * FROM t WHERE key BETWEEN 20 AND 60":
     1. descend to the leaf holding 20;
     2. walk the leaf chain, reading leaves sequentially.
     No backtracking, no internal-node data: a straight line of reads.
\end{verbatim}

\begin{verbatim}
   B vs B+ at a glance:

                 B-tree               B+ tree
   data          in every node        only in leaves
   internal nodes        keys + data + children   keys + children
   range scan    scatter, no chain    leaf chain = sorted sequence
   space reuse   duplicate keys in internal nodes   keys duplicated only
                               between levels (routing copies)
   best for      general purpose     databases, ordered range queries
\end{verbatim}

\subsection{B* Trees}

The B* tree tightens the B-tree's occupancy rule: nodes are kept at
least \emph{two-thirds} full instead of half-full, by redistributing
keys with a sibling before splitting, and by splitting one full node
into \emph{three} when both it and its sibling are full. The result is
fewer, denser nodes---fewer I/Os for the same data---at the cost of a
more complex insertion.

\begin{verbatim}
   B* tree: before splitting, borrow from a neighbor.

   order 5, node X full, sibling Y half-full:

        BEFORE                              AFTER (redistribute)
        [ 15 ]                              [ 17 ]
       /      \                            /       \
   [1|5|10|14] [17|20]              [1|5|10|14|15]   [17|20]
   (full)      (room)

   Only when BOTH are full does a B* split produce THREE nodes,
   each about 2/3 full, pushing two keys into the parent.

   cost/benefit:
     B-tree   min 1/2 full   simpler
     B* tree  min 2/3 full   33% fewer nodes, more complex inserts
\end{verbatim}

\begin{ruleofthumb}[When to use which tree]
Use a BST for in-memory data where simplicity wins and the data is
unpredictable. Use a B-tree when nodes map to blocks (disk, flash,
caches). Use a B+ tree when ordered range scans matter, as in nearly
every real database index. Use a B* tree when storage density
dominates and the extra insertion complexity is justified---they are
rare outside research and niche storage engines.
\end{ruleofthumb}

\begin{researchnote}[The B-tree family in the research record]
The B-tree has one of the longest and best-documented research
histories in computer science, and two papers in particular define
what you learned in this section. The first is the paper that
\emph{invented} the B-tree (Bayer and McCreight, 1972). The second is
the survey that catalogued its influence on databases (Comer, 1979).

\begin{verbatim}
   Extracts (condensed for the book):

       Bayer & McCreight (1972): the B-tree was designed so that
       searching a large ordered file requires only a few page reads,
       by keeping the tree balanced and each node the size of one
       storage page.

       Comer (1979), on the family's reach: the B-tree has become the
       standard method for organizing indexes in file systems and
       databases, and variants -- B+, B*, prefix B-trees -- were
       introduced to address specific needs.
\end{verbatim}

These extracts are paraphrases, not exact sentences; what matters is
the claim they establish: the wide, shallow, block-sized B-tree is not
a niche idea but \emph{the} industry-standard index structure, and the
B+ variant we implemented exists because ordered range scans are what
real queries need. When you meet \code{btree} indexes in every
production database, you are looking at a research lineage from 1972.

\paperref{D.~Comer, ``The Ubiquitous B-Tree,'' ACM Computing Surveys
11(2), 1979}{https://infolab.usc.edu/csci585/Spring2010/den_ar/indexing.pdf}

\paperref{R.~Bayer, E.~McCreight, ``Organization and Maintenance of
Large Ordered Indexes,'' Acta Informatica 1(3), 1972}{https://dl.acm.org/doi/10.1007/BF00288683}
\end{researchnote}

\section{Graphs}

\subsection{Graphs as Structs and Arrows}

A graph is a set of \emph{vertices} connected by \emph{edges}. Where a
tree is a graph with exactly one path between any two nodes and a
designated root, a graph drops every such restriction: cycles, multiple
edges, and unreachable regions are all allowed. Graphs model networks of
every kind---the internet, social graphs, road maps, dependency graphs,
circuit nets---and they are the most general structure in this chapter.

A graph is \emph{directed} if its edges have a direction (a one-way
street), \emph{undirected} if they do not (a two-way street), and
\emph{weighted} if each edge carries a cost. These are three
independent properties; all eight combinations exist.

\begin{verbatim}
   Undirected graph:                      Directed graph:

         0 -- 1                              0 --> 1
         |    |                              ^     |
         |    |                              |     v
         2 -- 3                              2 <-- 3

   edges are unordered pairs           edges are ordered pairs
   0-1 == 1-0                          0->1 is NOT 1->0

   Weighted (undirected):

            5
       0 -------- 1
       |          |
       3          2
       |          |
       4 -------- 3
            1

   each edge carries a number:         DAG (directed acyclic graph):
   the cost of traveling it.
                                          A --> B --> D
                                          |     ^
                                          v     |
                                          C --> E
                                          \
                                           v
                                           F
                                          (no cycle exists; topo-sortable)
\end{verbatim}

Two representations dominate, and they are the same trade-off we have
seen all chapter: pointer-honest locality versus constant-time
directness.

\subsection{Adjacency List}

Each vertex owns a linked list of its edges. The list is exactly the
linked list of this chapter; the graph is an array of such lists.

\begin{lstlisting}[caption={Adjacency list representation.}]
typedef struct graph_edge {
    int to;                  /* the vertex at the other end */
    int weight;              /* cost of the edge            */
    struct graph_edge *next; /* next edge leaving this vertex */
} graph_edge;

typedef struct graph {
    int n;                   /* number of vertices 0..n-1  */
    graph_edge **adj;        /* adj[v] = head of v's edge list */
} graph;
\end{lstlisting}

\begin{verbatim}
   The directed graph  0 -> 1 -> 2
                        ^    |
                        |    v
                        3 <- 4

   as an adjacency list (each vertex's outgoing edges):

     adj[0] -> (to=1) -----------> NULL
     adj[1] -> (to=2) -----------> NULL
     adj[2] -> NULL
     adj[3] -> (to=0) -----------> NULL
     adj[4] -> (to=1) -> (to=3) -> NULL

   Memory is O(V + E): one node per edge, one head per vertex.
   Asking "who are the neighbors of v?" is walking v's list.
   Asking "is u adjacent to v?" is walking u's list (or v's).
\end{verbatim}

Adding an edge pushes onto a vertex's list. For an undirected graph,
every edge must be added twice---once in each direction---and a
weighted edge stores its cost in the node.

\begin{lstlisting}[caption={Add a directed edge, and its undirected twin.}]
/* Add edge u -> v with the given weight. Returns 1 on success. */
int graph_add_edge(graph *g, int u, int v, int weight) {
    graph_edge *e;

    e = malloc(sizeof(graph_edge));
    if (e == NULL) {
        return 0;
    }
    e->to = v;
    e->weight = weight;

    e->next = g->adj[u];
    g->adj[u] = e;
    return 1;
}

/* Add an undirected edge: it is two directed edges. */
int graph_add_undirected(graph *g, int u, int v, int weight) {
    int r1;
    int r2;

    r1 = graph_add_edge(g, u, v, weight);
    r2 = graph_add_edge(g, v, u, weight);
    if (r1 == 0 || r2 == 0) {
        return 0;
    }
    return 1;
}
\end{lstlisting}

\begin{verbatim}
   Memory for one undirected edge u-v:

     adj[u] -> (to=v, w=w) -> ...
     adj[v] -> (to=u, w=w) -> ...

   Two directed edges model one undirected edge. The weight is stored
   on both sides, so both directions cost the same--as an undirected
   edge must.
\end{verbatim}

\subsection{Adjacency Matrix}

The alternative is a 2D matrix where entry \code{m[u][v]} is 1 (or the
weight) when the edge \code{u -> v} exists, and 0 otherwise. Any graph
problem that asks ``is this edge present?'' or needs all-pairs
information becomes a constant-time lookup.

\begin{lstlisting}[caption={Adjacency matrix representation.}]
typedef struct matrix_graph {
    int n;
    int **m;   /* m[u][v] = weight, or 0 when no edge u->v exists */
} matrix_graph;
\end{lstlisting}

\begin{verbatim}
   Directed graph  0 -> 1 -> 2
                    ^    |
                    |    v
                    3 <- 4

   as a matrix:

         0  1  2  3  4
       +----------------+
     0 | 0  1  0  0  0 |     m[0][1] = 1 : the edge 0->1 exists
     1 | 0  0  1  0  0 |     m[4][3] = 1 : the edge 4->3 exists
     2 | 0  0  0  0  0 |     every diagonal 0 : no self-loops
     3 | 1  0  0  0  0 |
     4 | 0  1  0  1  0 |
       +----------------+

   Memory is O(V^2) even for a sparse graph; each edge check is O(1).
   For a WEIGHTED matrix, a 0 is "no edge" and the value is the cost.
\end{verbatim}

\begin{verbatim}
   adjacency list vs adjacency matrix:

   concern              adjacency list      adjacency matrix
   ----------------     -----------------   ------------------
   memory               O(V + E)            O(V^2)
   edge check (u,v)?    O(degree) walk      O(1)
   iterate neighbors    O(degree), cached   O(V), must scan a row
   add edge             O(1)                O(1)
   best for             sparse, real nets   dense, small, or
                                              all-pairs algorithms
\end{verbatim}

\subsection{DAGs and Topological Order}

A DAG---directed acyclic graph---is a graph with no cycles. DAGs are
the shape of every dependency system: build systems, package managers,
compilation order, and spreadsheet cell evaluation all order a set of
jobs so that every job runs after its prerequisites. A \emph{topological
order} is an ordering of the vertices in which every edge points
forward. A graph admits one if and only if it is a DAG.

\begin{verbatim}
   A build system as a DAG (an arrow u -> v means "u must build
   before v"):

        lex -----------+
                       v
        yacc -----> make ---> main.o ---> myprog
                   ^                            ^
                   |                            |
                   +----------------------------+   (yacc also feeds myprog)

   Topological order must list each target after all of its inputs.
   Here: lex, yacc, make, main.o, myprog is valid; yacc before make
   is valid; make before lex is NOT (lex is a dependency of make).

   If a dependency graph has a cycle, no order exists: A needs B,
   B needs A, someone has to come first and cannot.
\end{verbatim}

Kahn's algorithm produces a topological order without recursion: count
every vertex's \emph{indegree} (incoming edges), then repeatedly emit a
vertex with zero indegree and delete it, decrementing the indegrees of
its successors. If we run out of zero-indegree vertices before emitting
everything, a cycle exists.

\begin{lstlisting}[caption={Topological sort by Kahn's algorithm.}]
#include <stdlib.h>

/* Write a topological order of g into order[]. Returns 1 on success,
   0 if g has a cycle. */
int graph_toposort(const graph *g, int *order) {
    int *indegree;
    int *queue;
    int head;
    int tail;
    int v;
    int emitted;
    const graph_edge *e;
    size_t i;

    indegree = calloc((size_t) g->n, sizeof(int));
    queue = malloc(sizeof(int) * (size_t) g->n);
    if (indegree == NULL || queue == NULL) {
        free(indegree);
        free(queue);
        return 0;
    }

    /* Pass 1: count incoming edges for every vertex. */
    for (i = 0; i < (size_t) g->n; i++) {
        e = g->adj[i];
        while (e != NULL) {
            indegree[e->to] = indegree[e->to] + 1;
            e = e->next;
        }
    }

    /* Seed the queue with the vertices that depend on nothing. */
    head = 0;
    tail = 0;
    for (i = 0; i < (size_t) g->n; i++) {
        if (indegree[i] == 0) {
            queue[tail] = (int) i;
            tail = tail + 1;
        }
    }

    /* Pass 2: emit zero-indegree vertices, one at a time. */
    emitted = 0;
    while (head < tail) {
        v = queue[head];
        head = head + 1;
        order[emitted] = v;
        emitted = emitted + 1;

        /* Removing v lowers the indegree of each successor. */
        e = g->adj[v];
        while (e != NULL) {
            indegree[e->to] = indegree[e->to] - 1;
            if (indegree[e->to] == 0) {
                queue[tail] = e->to;
                tail = tail + 1;
            }
            e = e->next;
        }
    }

    free(indegree);
    free(queue);

    if (emitted != g->n) {
        return 0;   /* some vertices never reached zero indegree: cycle */
    }
    return 1;
}
\end{lstlisting}

\begin{verbatim}
   Kahn's algorithm on the DAG above, step by step:

   indegrees:  lex=0 yacc=0 make=2 main.o=1 myprog=2
   (lex and yacc depend on nothing, so they are the "sources")

   queue: [lex, yacc]
   emit lex    ->  make indegree 2-1 = 1
   emit yacc   ->  make 1-1 = 0, enqueue make; myprog 2-1 = 1
   emit make   ->  main.o 1-1 = 0, enqueue main.o
   emit main.o ->  myprog 1-1 = 0, enqueue myprog
   emit myprog

   order: lex, yacc, make, main.o, myprog    -- every edge points forward.
\end{verbatim}

\subsection{Traversal: DFS and BFS}

Graphs are explored by two walks, and every serious graph algorithm
rides on one of them. \emph{Depth-first search} (DFS) dives down a
path as far as it goes before backtracking; \emph{breadth-first search}
(BFS) explores in expanding waves. Both mark visited vertices so they
never revisit, and both visit each vertex and edge once:
\code{O(V + E)}.

\begin{lstlisting}[caption={DFS, written with an explicit recursion.}]
/* Print the graph's vertices reachable from start, depth-first. */
void graph_dfs(const graph *g, int start, int *visited) {
    const graph_edge *e;

    visited[start] = 1;
    printf("%d ", start);

    e = g->adj[start];
    while (e != NULL) {
        if (visited[e->to] == 0) {
            graph_dfs(g, e->to, visited);
        }
        e = e->next;
    }
}
\end{lstlisting}

\begin{verbatim}
   DFS order from 0 on:                graph:
   visit 0, then 0's first neighbor 1,      0 - 1
   then 1's first unvisited neighbor 3,      | / |
   then 3's neighbor 2, done; backtrack.    2 - 3
                                           (undirected)

   order: 0, 1, 3, 2
   DFS goes as deep as it CAN before trying another neighbor of 0.
\end{verbatim}

BFS needs a queue; we reuse the array queue pattern from Kahn's
algorithm.

\begin{lstlisting}[caption={BFS, exploring in waves with a queue.}]
#include <stdlib.h>

/* Print the vertices reachable from start, breadth-first. */
void graph_bfs(const graph *g, int start, int *visited) {
    int *queue;
    int head;
    int tail;
    int v;
    const graph_edge *e;

    queue = malloc(sizeof(int) * (size_t) g->n);
    if (queue == NULL) {
        exit(1);
    }
    head = 0;
    tail = 0;

    visited[start] = 1;
    queue[tail] = start;
    tail = tail + 1;

    while (head < tail) {
        v = queue[head];
        head = head + 1;
        printf("%d ", v);

        e = g->adj[v];
        while (e != NULL) {
            if (visited[e->to] == 0) {
                visited[e->to] = 1;
                queue[tail] = e->to;
                tail = tail + 1;
            }
            e = e->next;
        }
    }

    free(queue);
}
\end{lstlisting}

\begin{verbatim}
   BFS order from 0 on the same graph:

   queue: [0]
   pop 0, enqueue 1, 2          queue: [1, 2]
   pop 1, enqueue 3             queue: [2, 3]
   pop 2 (3 already visited)    queue: [3]
   pop 3                        order: 0, 1, 2, 3

   BFS visits in distance order: everything one hop away, then two,
   then three. That is why BFS finds SHORTEST paths in an unweighted
   graph, while DFS is the natural tool for reachability and for
   detecting cycles.
\end{verbatim}

\begin{verbatim}
   DFS vs BFS at a glance:

                DFS                         BFS
   tool         recursion (or explicit      queue
                stack)
   order        deep before wide            distance order
   best for     reachability, cycles,       shortest paths, routing,
                topological, mazes          level order, social hops
   memory       O(depth)                    O(width)
   both run in  O(V + E), each vertex/edge visited once
\end{verbatim}

\section{Structures in Bytes}

Every container in this chapter is a struct, so every container has a
byte size. On a 64-bit machine (\code{int}=4, pointer=8):

\begin{verbatim}
   structure                node bytes   layout (offsets)
   ----------------------   ----------   ---------------------------
   singly linked list       16           value:0  next:8
   doubly linked list       24           value:0  prev:8  next:16
   circular list            16           value:0  next:8
   intrusive list_head      16           next:0   prev:8  (embedded)
   unrolled node            4 + 4B + 8   count:0 data:4 next:4B+8
   XOR list                 16           value:0  link:8
   skip list node           variable     value:0 level:4 forward[] at 8
   hash-table node          24           key:0 value:8 next:16
   B-tree node (order 4)    56           nkeys:0 keys:4 child:16 leaf:52
   graph edge               16           to:0  weight:4  next:8
\end{verbatim}

Two rows are worth reading twice. The B-tree node rounds up to 56
bytes because the \code{child[4]} pointer array must start at a
multiple of 8 (offset 16) and the trailing \code{leaf} int plus
padding fill the last 8. The skip-list node is variable because its
\code{forward[]} is a flexible array member: its true size is
\code{sizeof(skiplist\_node) + level * sizeof(pointer)}. These numbers
are what decide how many nodes fit in one 64-byte cache line---the
locality argument every data structure in this chapter is built on.

\section{Summary of the Menagerie}

Every structure in this chapter is the same two tools in a different
uniform: structs for shape, pointers for arrows, and a rule for who
owns each allocation. The cost table below is the one page you should
be able to redraw from memory.

\begin{verbatim}
   structure           search       insert     delete     notes
   ------------------  -----------  ---------  ---------  --------------------
   singly linked       O(n)         O(1)+      O(n)       1 ptr/node
   doubly linked       O(n)         O(1)+      O(1)+      2 ptrs/node
   circular            O(n)         O(1)+      O(1)+      ring, tail ptr
   intrusive           O(n)         O(1)+      O(1)+      links embedded
   unrolled            O(n)         O(n)       O(n)       blocks, cache
   XOR                 O(n)         O(1)+      O(1)+      1 ptr, needs prev
   skip list           O(log n)*    O(log n)*  O(log n)*  towers, sorted
   hash (chaining)     O(1)*        O(1)*      O(1)*      buckets+chains
   hash (open addr)    O(1)*        O(1)*      O(1)*      probing
   BST                 O(h)         O(h)       O(h)       h worst-case = n
   B-tree / B+/B*      O(log n)     O(log n)   O(log n)   wide, shallow
   graph adj list      O(deg)       O(1)       O(1)       O(V+E) memory
   graph adj matrix    O(1)         O(1)       O(1)       O(V^2) memory

   +  given a position/neighbor already found
   *  expected / probabilistic
\end{verbatim}

The four questions that opened the chapter are the checklist you apply
to every one of these---and to every structure you invent:

\begin{verbatim}
   SHAPE   : draw the struct.  How many pointers?  What points where?
   COST    : what does search/insert/delete actually touch?
   OWNER   : who mallocs, who frees, and what happens on the error path?
   CACHE   : are the bytes the algorithm wants in the same cache line?
\end{verbatim}

Answer those four honestly and you have designed a data structure;
code is the last, easy part. With this toolbox in hand we turn from the
shape of data to the services of the operating system: files, syscalls,
processes, signals, and the machinery every real C program sits on.





